\chapter{Hạn chế, kết luận và hướng phát triển}
\label{chap:limits}

Nhóm liệt kê các hạn chế một cách trung thực, để người đọc phân biệt phần đã chạy thật với phần còn ở mức demo hoặc thiết kế.

\section{Hạn chế hiện tại}

\begin{enumerate}
  \item \textbf{BR1--BR3 đã có đường dịch vụ, nhưng sandbox và provider vẫn là mức demo.} Dịch vụ thông báo và dịch vụ chấm bài chạy Kafka/Redis, API đọc mô hình Redis và chuyển bài nộp qua BFF; BR3 ghi schema riêng \code{grading.*}, còn projector sang LMS \code{public.*} là tùy chọn. Sandbox vẫn là tiến trình con và SQLite cục bộ, không phải ranh giới bảo mật; email/push cần thông tin xác thực provider thật để dùng ngoài local.
  \item \textbf{BR2 chưa có mô hình học máy.} Cả hai nhánh đặc trưng đã được hiện thực và phục vụ qua Redis (bảng lô được materialise, phiên chạy streaming), nhưng thứ hạng vẫn do luật quyết định; chưa có mô hình học từ dữ liệu hay bộ nhãn thật để so sánh.
  \item \textbf{Độ trễ streaming ở mức giây.} Spark Structured Streaming cho cảnh báo BR5 sau khoảng 10 giây và lớp phủ sau khoảng 7 giây, cao hơn mức dưới một giây mà Flink có thể đạt. Đây là đánh đổi có chủ đích (ADR 0009, 0014).
  \item \textbf{Embedding mang tính từ vựng.} Bộ nhúng băm không hiểu ngữ nghĩa (``join'' không khớp ``merge''). Đã có cổng \code{Embedder} và đường reindex để thay bằng mô hình thật.
  \item \textbf{Lớp phủ không mang phiên bản nội dung.} Vào ngày một video được quay lại, lớp phủ trộn hai phiên bản; vì vậy nó chỉ được hiển thị cho phiên bản mới nhất.
  \item \textbf{Lịch sử BR5 trong kho dữ liệu đến từ dữ liệu tổng hợp.} Runtime BR4/BR5 lưu điểm và cảnh báo trong PostgreSQL; cần phát chúng lên \code{risk.learner_score.v1} và \code{risk.dropout_alert.v1} (hoặc thêm bảng vào connector CDC) để kho dữ liệu nhận từ đường thật.
  \item \textbf{Hiệu năng BR4.} Mỗi lần sửa dữ liệu làm mới toàn bộ bản tổng hợp bị ảnh hưởng và dùng phép nối JSONB; với 100.000 lượt chấm, một lần sửa tốn khoảng 4 giây. \todofix{xác nhận lại BR4}
  \item \textbf{Vận hành đơn nút, bảo mật cục bộ.} Trino, OpenSearch, Kafka chạy một nút; plugin bảo mật của OpenSearch tắt; token tĩnh thay cho hệ thống định danh của trường. Chưa có giám sát (độ trễ consumer, thời gian checkpoint, SLA độ tươi).
  \item \textbf{Dữ liệu tổng hợp, không nhãn thật.} Không có nhãn bỏ học thật nên không đánh giá được độ chính xác của cảnh báo; các số đo hiệu năng là quan sát trên một máy.
  \item \textbf{Lệch dữ liệu mẫu.} Seed của lakehouse khai báo khoảng 18 khoá học và video không có trong cơ sở dữ liệu LMS mẫu, nên job reindex báo các mục này là lỗi.
\end{enumerate}

\section{Hướng phát triển}

\begin{itemize}
  \item Đưa projector BR3 vào CI/e2e và thử bật đồng bộ LMS trên một bản sao dữ liệu, kèm quy trình re-baseline báo cáo.
  \item Thử một mô hình gợi ý học từ dữ liệu cho BR2, so sánh với luật hiện tại; hai nhánh đặc trưng đã có trong Redis.
  \item Chuyển các job cần độ trễ dưới một giây sang Flink khi nền tảng lên Flink 2.x và connector Delta Kernel ổn định.
  \item Thay bộ nhúng băm bằng mô hình sentence-embedding, chạy một lần reindex bắt buộc.
  \item Chuyển các bảng snapshot (\code{fct_enrollment}, \code{fct_lesson_progress}) sang MERGE dựa trên Change Data Feed khi khối lượng đủ lớn; với dashboard dưới giây ở quy mô lớn, gắn StarRocks vào cùng Hive Metastore.
  \item Triển khai lên Kubernetes với nhiều nút, TLS và xác thực tập trung; giám sát bằng Prometheus cùng Grafana; bổ sung quy trình xoay khoá bút danh.
  \item Thu thập nhãn kết quả thật để hiệu chỉnh và đánh giá điểm rủi ro BR5.
\end{itemize}

\section{Kết luận}
\label{chap:conclusion}

Báo cáo đã trình bày DE\_261, một nền tảng kỹ thuật dữ liệu biến bốn nguồn dữ liệu rời rạc của một hệ thống giáo dục trực tuyến thành các kho dữ liệu phục vụ tám yêu cầu nghiệp vụ.

Về kiến trúc, nền tảng bám sát kiến trúc tham chiếu của đề bài: Kafka và Debezium cho tiếp nhận, Spark cho xử lý, Delta Lake cho lakehouse, Redis cho phục vụ độ trễ thấp. Lớp gold được mô hình hoá thành kho dữ liệu Kimball và phục vụ qua Trino. Các điểm điều chỉnh, quan trọng nhất là dùng Spark Structured Streaming thay Flink cho BR5 và các lớp phủ, đều được ghi lại bằng ADR kèm lý do và hệ quả. Tổ chức mã theo DDD và kiến trúc lục giác cho phép cùng một logic nghiệp vụ chạy trên nhiều engine và được kiểm thử độc lập.

Về hiện thực, BR4--BR8 chạy trên hạ tầng thật với các tính chất mà một hệ thống dữ liệu cần có: xử lý theo thời gian sự kiện, chịu được dữ liệu trễ, trùng lặp và sai thứ tự, chạy lại idempotent, khử định danh ở ranh giới silver, k-anonymity cho báo cáo tổ chức, và cổng kiểm tra chất lượng sau mỗi lớp. BR1--BR3 có logic nghiệp vụ hoàn chỉnh và chạy trong runtime demo. Một giao diện web mới đưa kết quả đến đúng từng nhóm người dùng.

Về đánh giá, hệ thống có hơn 360 test tự động ở nhiều tầng và các kịch bản kiểm chứng nghiệp vụ được quan sát đến tận giao diện. Các số đo cho thấy độ trễ phù hợp với từng yêu cầu: cảnh báo bỏ học khoảng 10 giây sau thay đổi nguồn, nội dung mới tìm kiếm được sau khoảng nửa giây, mart được phục vụ trong vài mili-giây khi có cache.

Bài học lớn nhất của nhóm là giá trị của việc kiểm chứng bằng dữ liệu và số đo. Nhiều lỗi quan trọng chỉ lộ ra khi chạy hệ thống thật và đo đạc. Việc ghi rõ phần nào đã chạy, phần nào còn là thiết kế giúp nhóm và người đọc đánh giá đúng mức độ hoàn thiện, và định hướng cho các bước phát triển tiếp theo.
